\chapter{Model Hibrida: Fisika dan Data Bersama}
\label{ch:hibrida}

Di antara semua pendekatan di Bagian V, model hibrida adalah yang paling dekat dengan cara kerja
insinyur sehari-hari. Kita tidak membuang solver yang sudah teruji puluhan tahun. Kita membiarkannya
mengerjakan apa yang ia kuasai, lalu memakai data untuk memperbaiki bagian yang ia kerjakan dengan
buruk. Bab ini dirangkai dari kuliah metode numerik fluida, yang membahas model penutup turbulensi
beserta kelemahannya, dan dari bacaan selama tugas akhir.

\section{Di mana solver lemah}

\Cref{ch:numerik} memperlihatkan bahwa simulasi praktis selalu memodelkan sesuatu yang tidak bisa
diselesaikan: tegangan Reynolds di RANS, pusaran sub-grid di LES, gaya antara fluida dan partikel di
aliran multifase, atau perpindahan panas di dekat dinding yang tidak cukup halus grid-nya. Model-model ini
dirumuskan dengan asumsi yang disederhanakan dan dikalibrasi pada aliran sederhana. Ketika dipakai pada
aliran yang rumit, misalnya dengan pemisahan aliran, kelengkungan kuat, atau putaran, asumsinya runtuh.
Data dari simulasi yang lebih teliti, seperti DNS dan LES, atau dari percobaan, bisa dipakai untuk
mempelajari model yang lebih baik \citep{duraisamy2019}.

\section{Model penutup yang dipelajari}

Pendekatan paling langsung adalah melatih jaringan untuk memprediksi besaran yang hilang dari besaran
yang tersedia. Untuk RANS, targetnya tegangan Reynolds, dan inputnya gradien kecepatan rata-rata. Tetapi
jaringan sembarang tidak menjamin sifat fisika yang paling dasar: hasilnya harus sama tidak peduli
bagaimana sistem koordinat diputar. \citet{ling2016} membangun sifat itu ke dalam arsitektur. Mereka
memakai hasil klasik bahwa tensor anisotropi Reynolds $\bm b$, dalam model viskositas turbulen umum,
bisa ditulis sebagai kombinasi dari sepuluh tensor basis $\bm T^{(n)}$ yang disusun dari tensor laju
regangan $\bm S$ dan tensor rotasi $\bm\Omega$ yang dinormalisasi \citep{pope2000}:
\begin{equation}
\bm b=\sum_{n=1}^{10}g_n(\lambda_1,\dots,\lambda_5)\,\bm T^{(n)}.
\end{equation}
Koefisien $g_n$ hanya bergantung pada lima besaran invarian $\lambda_i$, seperti
$\mathrm{tr}(\bm S^2)$ dan $\mathrm{tr}(\bm\Omega^2)$, yang nilainya tidak berubah ketika koordinat
diputar. Tensor basis sendiri berputar bersama koordinat. Jadi kalau jaringan hanya mempelajari $g_n$
dari invarian, hasil $\bm b$ dijamin berputar dengan benar, apa pun bobot jaringannya. Model viskositas
turbulen linear dari \cref{ch:numerik} adalah kasus khusus dengan hanya suku pertama, $\bm T^{(1)}=\bm S$,
dan koefisien tetap. Inilah contoh nyata inductive bias dari \cref{ch:peta}: simetri tidak dipelajari dari
data, tetapi dibangun ke dalam bentuk modelnya.

\section{A priori dan a posteriori}

Model penutup yang dipelajari sering tampak baik dalam uji \emph{a priori}: ambil medan DNS,
hitung input dan target, lalu ukur galat prediksi. Begitu model dimasukkan ke dalam solver dan simulasi
dijalankan, yaitu uji \emph{a posteriori}, hasilnya bisa mengecewakan atau bahkan meledak. Penyebabnya
sama dengan masalah rollout di \cref{ch:operator}. Selama pelatihan, model melihat input dari DNS.
Selama simulasi, ia melihat input dari medan yang dihasilkan solver bersama model itu sendiri, yang
distribusinya bergeser. Galat kecil mengubah medan, medan yang berubah mengubah input model, dan
lingkarannya bisa memperbesar galat.

Obatnya adalah melatih model di dalam lingkar yang sama dengan tempat ia akan dipakai. Pendekatan
\istilah{solver-in-the-loop} \citep{um2020} menulis satu langkah simulasi hibrida sebagai
\begin{equation}
\vu_{n+1}=\mathcal S(\vu_n)+\mathcal C_\theta\big(\mathcal S(\vu_n)\big),
\end{equation}
dengan $\mathcal S$ solver kasar dan $\mathcal C_\theta$ koreksi yang dipelajari. Kalau solver ditulis
dalam kerangka yang bisa diturunkan, gradien loss bisa dibawa mundur melewati beberapa langkah solver
sekaligus, sehingga koreksi belajar memperbaiki galat yang muncul dari interaksinya sendiri dengan solver.
Gagasan solver yang bisa diturunkan, \istilah{differentiable physics}, dibahas lengkap dalam buku
\citet{thuerey2021}. \citet{kochkov2021} memakai pendekatan serupa untuk turbulensi dua dimensi: jaringan
mempelajari cara menginterpolasi nilai di sisi sel pada grid yang jauh lebih kasar, dan simulasi yang
dihasilkan setara akurasinya dengan solver klasik pada grid berkali-kali lebih halus, dengan percepatan
puluhan kali.

\section{Ketidakpastian dan banyak tingkat ketelitian}

Model hibrida juga cocok untuk memanfaatkan data dari sumber yang berbeda-beda ketelitian dan harganya.
Simulasi kasar murah dan banyak, simulasi teliti mahal dan sedikit, percobaan paling dipercaya tetapi
paling jarang. Model \istilah{multi-fidelity} mempelajari koreksi dari tingkat murah ke tingkat mahal, yang
biasanya jauh lebih sederhana untuk dipelajari daripada solusi lengkapnya. Ketidakpastian prediksi bisa
diperkirakan dengan ensemble atau model probabilistik (\cref{ch:arsitektur}), dan ini penting karena model
hibrida sering dipakai justru di luar daerah data latihnya.

Asimilasi data dari \cref{ch:prob} adalah bentuk hibrida yang paling tua. Solver memberi ramalan, data
memberi koreksi, dan keduanya digabung dengan bobot sesuai ketidakpastian masing-masing. Ketika model
penutupnya sendiri tidak pasti, parameter model bisa ikut ditaksir dalam proses asimilasi, misalnya
dengan menambahkannya ke vektor keadaan Kalman filter ensemble.

\section{Apa yang dipelajari, dan apa yang tidak}

Pola umum model hibrida bisa diringkas dalam satu kalimat: biarkan fisika menangani apa yang sudah
diketahui pasti, terutama hukum kekekalan, dan biarkan data menangani apa yang dimodelkan dengan buruk.
Pilihan ini punya keuntungan praktis yang besar. Karena sebagian besar dinamika dihitung oleh solver
yang konservatif, model hibrida jauh lebih jarang melanggar kekekalan massa atau energi daripada model
pengganti murni. Karena koreksinya kecil, ia butuh lebih sedikit data. Dan karena strukturnya jelas,
kegagalannya lebih mudah didiagnosis. Ulasan \citet{vinuesa2022} dan \citet{brunton2020} menempatkan
pendekatan ini sebagai salah satu jalur yang paling menjanjikan untuk mekanika fluida komputasi.

\sumberbab{Bab ini dirangkai dari kuliah Numerical Methods in Fluid Mechanics (model penutup turbulensi)
dan bacaan selama tugas akhir. Rujukan utama: \citet{duraisamy2019}, \citet{ling2016}, \citet{um2020},
\citet{kochkov2021}, dan \citet{thuerey2021}.}
